% !!!!!!!!!!!!!!!!!!!!!!!!!!!!!!!!!!!!!!!!!!!!!!!!!!!!!!!!!!!!!!!!!!!!!!!!!!!!!!!!
% !!!!              --- IMPORTANT FOR OVERLEAF USERS ---                      !!!!
% !!!!  You must make the following changes to the default settings for this  !!!!
% !!!!  to work                                                               !!!!
% !!!!  Go to File -> Settings -> Compiler and change in the drop down menus: !!!!
% !!!!    -  Compiler to XeLaTeX                                              !!!!
% !!!!    -  Main document to main.tex                                        !!!!
% !!!!!!!!!!!!!!!!!!!!!!!!!!!!!!!!!!!!!!!!!!!!!!!!!!!!!!!!!!!!!!!!!!!!!!!!!!!!!!!!
%-----------------------------  README  ------------------------------------------
% University of Leeds thesis template - Erin Raif, 2025 (hopefully)
% Inspired by templates of Sam Greenwood (2020) and Jonathan Coney (2024).
% **Tested** pull requests to improve this are very welcome!
%
% Compatible with University of Leeds guidelines (included below), and designed
% for use when submitting a standard thesis, rather than by publication.
% I recommend that those looking for template for a thesis by publication
% use the Sam Greenwood thesis template that can be found at
% https://github.com/sam-greenwood/SOEE-thesis-template
%
% PLEASE NOTE: If you click recompile while viewing this preamble, Overleaf
% will break! This is an overleaf bug. You can only recompile while viewing
% a file with content (e.g. main.tex; 1_intro.tex). If you're still struggling
% to make it work, check that File -> Settings -> Compiler -> Main Document is
% main.tex. For offline compilation, read full notes (which you should anyway).
%
% ------------ Notes on this template: ------------------------------------------
% STRUCTURE - each piece of writing is contained in a separate .tex file.
% This allows faster compilation by only selecting relevant files using the
% \includeonly command on line 17 of main.tex. 
%
% LINKING and CROSS-REFERENCING - this package is designed to use
% the \cref command rather than the \ref demand, automatically creating
% the e.g. "Figure" or "Table" part of the reference, which is clickable.
%
% FONT - Open Sans font is used here and .ttf files are supplied. There are
% two warnings since there are one or two unusal characters that don't have
% a monospace bold equivalent, these are safely ignored (and can be turned
% off with the silence pacakges). Alternatively, remove reference to fontspec,
% \setmainfont and \sffamily to revert to Computer Modern.
% Open Sans is a open-source humanist font designed by Steve Matteson.
%
% REFERENCING - The University of Leeds only requires you to use an
% "established referencing standard" in your field. Since I studied
% Atmospheric Physics, I have chosen the Copernicus style used by the
% Atmospheric Chemistry and Physics Journal, since the Leeds Harvard style
% is grim.
%
% PDF METADATA - You will need to change this in the preamble to include
% your own name, title etc. I could have automated this and the front page.
% I haven't.
%
% CUSTOM MATHS/HYPENATION/SHORTHANDS - enter these in custom_defs.tex
%    for ease and input after maths packages have been loaded
%
% ABBREVIATIONS - please see detailed comment in preamble to choose between
%    two possible setups for abbreviations. You need to make this decision
%    early, I'm afraid.
% 
% --- Complilation instructions: ---
%   - Tested to work on Overleaf with XeLaTeX (as of February 2026)
% TODO DO THIS OFFLINE, NOTE USING XeLaTeX NOT pdfLaTeX if using font changes here
%   - Otherwise to compile yourself:
%     pdflatex 0_thesis.tex
%     biber 0_thesis                                         
%.    makeindex 0_thesis.nlo -s nomencl.ist -o 0_thsis.nls    
%     pdflatex thesis.tex
% PDFLATEX WILL BREAK WHEN USING THE CUSTOM FONT!
%
% --- University of Leeds Thesis Guidance (Aug 2024) ---
% Please note that you should check the guidance has not changed since I
% started my thesis. To be honest, I will probably forget to check that
% the guidance hasn't changed by the time I submit it. I've only included
% the things that actually affect what needs to go into the LaTeX template,
% you should check for rules on the actual thesis content!
%
% Font - clear, standard font of 11-12pt size
% Text in double or 1.5 line spacing
% Margins at binding edge not less than 40mm
% Margins at other edges not less than 20mm
% Pages numbered consecutively, inclduing images or diagrams
% Page numbers on every page, preferably top centre (top-corner here)
% Title page must have thesis full title, your full name, submission status,
%     The University of Leeds, your School, Month and Year of Exam
%     note the date is not the date of submission!
% Second page must have intellectual property statements
% Page 3 - acknowledgements with at least those who have assisted work
%     as well as optional personal acknowledgements.
% Abstract of max. 300 words on page 4 (or 5 since we're using correct
%   binding)
% Table of Contents and further content lists on page immediately
%     after abstract page.
% Ethics approval in appendix, list of abbreviations somewhere
% References using an established standard in your subject area
%
% With additional thanks to preceding students (now doctors) whose work
% the Greenwood template was based on: M.C. Garthwaite, D.P.S. Baekart
% and M.E. Gaddes
%---------------------------------------------------------------------------------

\documentclass[titlepage,twoside,onecolumn,a4paper,11pt]{report}

%%%%%%%%%%%%%%%%%%%%%% Helpful LaTeX pacakges for macros %%%%%%%%%%%%%%%%%%%%%%%%%%%%%%%%%%%%%
\usepackage{ifthen}
\usepackage{comment}
\usepackage{silence}                        % get rid of specific annoying errors or warnings
                                            % using WarningFilter (example below)
                                            % not recommended, but I'm not your dad
                                            % and I definitely didn't use the filters below to
                                            % remove warnings for safe choices
% \WarningFilter{latexfont}{Font shape `TS1/OpenSans}
% \WarningFilter{latexfont}{Some font shapes were not available, defaults substituted}
%%%%%%%%%%%%%%%%%%%%%% Fonts, symbols and text %%%%%%%%%%%%%%%%%%%%%%%%%%%%%%%%%%%%%%%%%%%%%%%
% If you want to revert to traditional LaTeX Computer Modern font, just comment the font setting bits
% out and the switch at the end
\usepackage[T1]{fontenc}                    % Allow fonts with foreign charaters
\usepackage{fontspec}                       % Use different font from Computer/Latin Modern
\usepackage{textcomp}                       % Avoid warnings from gensymb

% Setting custom font - ADVANCED note. I have opted for a sans-serif font. To implement this
% I renew the familydefault command so that sans-serif is the default. This means you can
% still switch to a roman (seriffed) font if need be.
% I then use \setsansfont to set my custom font, which is OpenSans. If you want a Roman
% font of your own choosing, you should remove the renewal of the familydefault command below
% and change \setsansfont to \setmainfont.
\renewcommand{\familydefault}{\sfdefault}
\setsansfont{OpenSans}[                     % Use OpenSans font
    Path=./OpenSans/,                       % Note this is switched on at very end of preamble
    Scale=0.9,
    Extension = .ttf,
    UprightFont=*-Regular,                  % alternatively, "Medium" is provided
    BoldFont=*-SemiBold,                    % alternatively, "Bold" is provided
    ItalicFont=*-LightItalic,               % alternatively, "Italic" is provided
    BoldItalicFont=*-SemiBoldItalic         % alternatively, "BoldItalic" is provided
    ]
\usepackage[dvipsnames]{xcolor}             % Colour fonts, with extra named font options
\usepackage{textgreek}                      % Greek letters in text
\usepackage{gensymb}                        % Additional "generic symbols" like degrees
\newcommand{\textapprox}{\raisebox{0.5ex}{\texttildelow}} % in-line tilde in text
\usepackage[UKenglish]{babel}               % Ensures correct hyphenation
% Note for advanced users:
% Open Sans is a modern font for unicode, so you should be able to remove textcomp and gensymb
% and avoid the warnings - however, this requires a little more overhead to create symbols
% like degrees, so I've left them in. If you're trying to add something particularly complex,
% try removing these as they can interfere with some other packages.
%%%%%%%%%%%%%%%%%%%%%% Optional dark mode for night owls %%%%%%%%%%%%%%%%%%%%%%%%%%%%%%%%%%%%%
%\pagecolor[rgb]{0.075,0.075,0.075}         % "dark mode" background when uncommented
%\color[rgb]{0.9,0.9,0.9}                   % "dark mode" font when uncommented
                                            % note this does not change images/non-default 
                                            % colours produced with \textcolor
%%%%%%%%%%%%%%%%%%%%%% Page layout %%%%%%%%%%%%%%%%%%%%%%%%%%%%%%%%%%%%%%%%%%%%%%%%%%%%%%%%%%%
\usepackage{a4}                             % TODO: check if redundant
\usepackage{geometry}                       % Package to setup page margins
\geometry{a4paper,inner=40mm,outer=25mm,top=30mm,bottom=30mm}
\usepackage{multicol}                       % multicolumn text where appropriate (not used here
                                            % but you might want it to save reference space)
\usepackage{verbatim}                       % allows font without formatting
\usepackage{pdflscape}                      % for landscape pages in the appendix
\usepackage{setspace}                       % choose spacing
\onehalfspacing                             % my preference - 1.5x linespacing
%\doublespacing                             % if you prefer 2x linespacing
\usepackage{xurl}                           % neat URL format
\usepackage[left,mathlines]{lineno}         % Places line numebrs in margin
\usepackage{fancyhdr} 						% for making headers and footers
\pagestyle{fancy}                           % 'fancy' will place header and footers. 'plain' will not
\renewcommand{\chaptermark}[1]{\markboth{\chaptername\ \thechapter:\ #1}{}} % sets format of Chapter mark (left pages)
\renewcommand{\sectionmark}[1]{\markright{Section \thesection:\ #1}}        % sets format of section mark (right pages)
\renewcommand{\headrulewidth}{0.5pt}        % sets thickness of header rule
\renewcommand{\footrulewidth}{0pt}          % sets thickness of footer rule
\fancyhead[RO,LE]{\thepage} 				% options R - Right-hand side L - Left-hand side O - odd page E - even page
\fancyfoot{}
\setlength\headheight{13.6pt}			    % if chapter titles stay on one line in the headers
%\setlength\headheight{26.0pt}              % if chapter titles need to go onto two lines in the header
                                            % (note you can avoid this by using \chaptermark to have a shorter header)
\fancyhfoffset[EL,OR]{0mm}
%%%%%%%%%%%%%%%%%%%%%% Figures %%%%%%%%%%%%%%%%%%%%%%%%%%%%%%%%%%%%%%%%%%%%%%%%%%%%%%%%%%%%%%
\usepackage{graphicx}                       % Effectively default package for graphics.
\usepackage{placeins}                       % for more comfortable float placement including \FloatBarrier
\usepackage{wrapfig}                        % figures to side of page with text wrapping around them
\usepackage{rotating}                       % allow rotated floats (useful for sideways tables)
\usepackage[font={small},labelfont=bf,singlelinecheck=false]{caption} % caption with separate properties
\renewcommand{\figurename}{Figure}          % Sets figures to be presented as Fig. N
\graphicspath{ {./images/} }                % change as appropriate, allows you to use filenames only
\renewcommand{\dblfloatpagefraction}{.66}   % place floats on separate page more appropriately by
                                            % changing the size a figure has to be
                                            % (relative to the page size) to get its own page
%%%%%%%%%%%%%%%%%%%%%% Tables %%%%%%%%%%%%%%%%%%%%%%%%%%%%%%%%%%%%%%%%%%%%%%%%%%%%%%%%%%%%%%%
% Note these choices are compatible with LaTeX tables generated using 
% https://www.tablesgenerator.com/
% though you will want to swap tabular for tabularx in the table generator output should you
% want to avail of column types L and x that have been added below or take advantage of
% any of the other useful features of tabularx.
\usepackage{tabularx}                               % Perhaps a more modern choice is booktabs but this is what I'm used to
\usepackage{array}                                  % Definition of table sizes allowed
\newcolumntype{x}{>{\centering}X}                   % Neat centre-aligned columns in tabularx tables
\newcolumntype{L}{>{\raggedright\arraybackslash}X}  % Neat left-aligned columns in tabularx tables
\usepackage{multirow}                               % multirow boxes in tables where appropriate
\usepackage{xltabular}
%%%%%%%%%%%%%%%%%%%%%% Mathematics %%%%%%%%%%%%%%%%%%%%%%%%%%%%%%%%%%%%%%%%%%%%%%%%%%%%%%%%%%
\usepackage{amsmath}
\usepackage{bm}
\usepackage[separate-uncertainty=true]{siunitx}     % Consistent SI units. Use \qty, \unit and \qtyrange commands.
                                                    % e.g. \qty{20}{\per\centi\meter\cubed} is 20cm^-3 printed nice
% NOTE FOR ADVANCED USERS:
% The next line tells siunitx to use the text font for numbers throughout, even if they're in equations.
% This means writing, say, $<\qty{0.5}{\metre}$ in a sentence looks in keeping with the rest of the text.
% This isn't always what you want, so if you need units in a proper equation, use \qty[mode=math].
\sisetup{detect-weight,detect-shape,mode=text}      

%%%%%%%%%%%%%%%%%%%%%% Hyperlinks and PDF Metadata Setup %%%%%%%%%%%%%%%%%%%%%%%%%%%%%%%%%%%%
\usepackage[pdfpagelabels,hidelinks]{hyperref}      % Use in-document hyperlinks
\usepackage[capitalize,nameinlink]{cleveref}        % Use clever references to get names linking and autopopulation of type
                                                    % NOTE: Weird quirk - cleveref struggles to identify
                                                    % tables and figures from the main text if you do not
                                                    % indent within your figure/table environment. Very strange...
                                                    
\crefformat{figure}{#2Figure~#1#3}                  % Set figure reference in text to be Figure rather than Fig.
\hypersetup{
    plainpages=false,                               % Resets counter in main matter (Problems of links to early pages)
    %bookmarks=true,                                % show bookmarks bar?
    unicode=true,                                   % non-ASCII characters in Acrobat's bookmarks
    pdftoolbar=true,                                % show Acrobat's toolbar?
    pdfmenubar=true,                                % show Acrobat's menu?
    pdffitwindow=true,                              % page fit to window when opened
    pdfstartview={FitV},
    pdftitle={The formation of ice in clouds and the impact on climate.},           % title
    pdfauthor={Erin N. Raif},                                                       % author
    pdfsubject={PhD Thesis, School of Earth and Environment, University of Leeds},  % subject of the document
    pdfkeywords={},
    pdfdisplaydoctitle=true
}
\usepackage{doi}
\usepackage[nottoc]{tocbibind}      % Bibliography in table of contents
%%%%%%%%%%%%%%%%%%%%%% Appendices %%%%%%%%%%%%%%%%%%%%%%%%%%%%%%%%%%%%%%%%%%%%%%%%%%%%%%%%%%%
\usepackage{appendix}                       % Allows appendix environments
%%%%%%%%%%%%%%%%%%%%%% List of abbreviations (acronyms) %%%%%%%%%%%%%%%%%%%%%%%%%%%%%%%%%%%%%
% There are two options here and you should decide which to use BEFORE YOU START.
% Option 1: Abbreviations link to the list of abbreviations. Use the \ac command whenever you
%           write an acronym in the text, e.g. \ac{INP} to create the list. Pros and cons:
%           + will print the full name of the abbreviation on the first instance
%               e.g. ice-nucleating particle (INP) but not on others.
%           + links to the list every time, which is helpful
%           + you should (?) be warned if you call an undefined abbreviation
%           - extra overhead when writing to remember to do this every time
%           - sometimes stylistically you want to clarify what an acronym is midway through
%           - lots of hyperlinks can be irritating for your examiner if they read on tablet
% Option 2: Low-effort acronym implementation. Manually create the list of abbreviations. No
%           need to type \ac all the time but you do need to remember to add them as you go.
%\usepackage[printonlyused]{acronym}  % <--- Option 1
\usepackage[nohyperlinks]{acronym}    % <--- Option 2 (I used this)
%%%%%%%%%%%%%%%%%%%%%% Nomenclature (a.k.a. list of symbols) %%%%%%%%%%%%%%%%%%%%%%%%%%%%%%%%
% If you want to have groups in your nomenclature, you will need to use the etoolbox package
% and follow the instructions here. https://www.overleaf.com/learn/latex/Nomenclatures
% The Greenwood thesis template linked above has an example of implementing this
\usepackage{nomencl}                        
\makenomenclature                           
\setlength\nomlabelwidth{1.5cm} 
%%%%%%%%%%%%%%%%%%%%%% Bibliography %%%%%%%%%%%%%%%%%%%%%%%%%%%%%%%%%%%%%%%%%%%%%%%%%%%%%%%%%
% This uses the old-fashioned BibTeX way of doing references. This works nicely with .bst
% files so you can pick your favourite (mine is Copernicus) and put straight in. However,
% it does not do references by chapter - see the Greenwood thesis template which uses biber
% rather than BibTeX to implement this
\usepackage[comma]{natbib}          % Author-date referencing style
\setcitestyle{citesep={;}}          % Separate in-line refreences with semi-colons        
% NOTE 2026: Copernicus now has an updated .bst style. This template has been tested with
% the old version. There's no reason why the new version shouldn't work, so I cautiously
% recommend you upgrade if you encounter any anonying bibliography issues.
%%%%%%%%%%%%%%%%%%%%%% Thesis-specific commands (my/your hacks)%%%%%%%%%%%%%%%%%%%%%%%%%%%%%%
%------- Custom definitions -------%
% Maths symbols and shorthands are defined for your convenience in a separate file
% I also include words which the language package struggles to hyphenate over two lines
%%%%%%%%%%%%%%%%% Definitions of custom LaTeX macros %%%%%%%%%%%%%%
% Derivatives
\newcommand*\diff{\mathop{}\!\mathrm{d}} 
\newcommand*\Diff[1]{\mathop{}\!\mathrm{d^#1}} 
% CASIM hydrometeors
\newcommand\liquid{\mathrm{w}}
\newcommand\water{\mathrm{w}}
\newcommand\rain{\mathrm{r}}
\newcommand\ice{\mathrm{i}}
\newcommand\snow{\mathrm{s}}
\newcommand\graupel{\mathrm{g}}
\newcommand\vapour{\mathrm{v}}
% Types of INP concentration
\newcommand{\inpconc}{N_\mathrm{INP}}
\newcommand{\inpcontact}{N_{\text{INP, con.}}}
\newcommand{\inpimmersion}{N_{\text{INP, imm.}}}

% Hyphenation that babel didn't know how to do.
\hyphenation{super-satur-ation}
%------- Subsubsubsections -------%
% Create subsubsubsection (e.g. 3.1.1.1) that appears in the Table of Contents.
% Call by using the \paragraph command. God help the poor reader if you need it!
\usepackage{titlesec}                       
\setcounter{secnumdepth}{4}
\titleformat{\paragraph}
{\normalfont\normalsize\bfseries}{\theparagraph}{1em}{}
\titlespacing*{\paragraph}
{0pt}{3.25ex plus 1ex minus .2ex}{1.5ex plus .2ex}
%------- Chapters start on right-hand page -------%
% Define "clear double page". First end the current page and start a new one, dumping floats.
% Then check for two-sided formatting. If so, and the current page number is odd, continue
% on the current page. If the current page no. is even, make an empty line on an empty page,
% then make a new page (unless two-column, in which case make a new empty column first).
\makeatletter
\def\cleardoublepage{\clearpage\if@twoside \ifodd\c@page\else%
\hbox{}%
\thispagestyle{empty}
\newpage%
\if@twocolumn\hbox{}\newpage\fi\fi\fi}
\makeatother
%-------  Include thesis chapter -------%
% Create command to include thesis chapters. Takes one argument: chapter filename.
% Note unlike for precursors, the chapter env is set up INSIDE the chapter.tex file.
% See main.tex for discussion of how this appears to create unexpected blank pages! 
\newcommand{\includethesischapter}[1]{
  \include{#1}
  \cleardoublepage
  \fancyhead[RE]{\leftmark}
  \fancyhead[LO]{\rightmark}
}
%%%%%%%%%%%%%%%%%%%%%% Font %%%%%%%%%%%%%%%%%%%%%%%%%%%%%%%%%%%%%%%%%%%%%%%%%%%%%%%%%%%%%%%%%
% If you don't want to use OpenSans but do want a sans-serif font, you can use the default
% built-in one with zero faff by uncommenting this line. Make sure you remove all the guff
% that I've added if you do!
%\sffamily
